% !TeX root

\documentclass[12pt,a4paper]{article}
% =============================================================
% Preamble entry point. Each file in config/ handles one concern.
% Order matters: do not reorder without checking package loading.
% =============================================================
% FONTS & LANGUAGE
% Font encoding: T1 for proper accents and hyphenation
\usepackage[T1]{fontenc}
% Language: British English (matches A4 paper).
% For American English, change to main=american.
\usepackage[main=british]{babel}
% Fonts: Times-like text and math (newtx)
\usepackage{newtxtext,newtxmath}
% Quotes: style follows the babel language; write \enquote{...} in the text.
% Must be loaded after babel and before biblatex.
\usepackage[autostyle=true]{csquotes}
% Typography: better justification and hyphenation
\usepackage{microtype}

% Note: no inputenc needed, UTF-8 is the default in modern LaTeX.
          % fonts, babel (English), csquotes
% PAGE GEOMETRY & TITLE POSITION
% 1-inch margins on all sides
\usepackage[a4paper,margin=1in]{geometry}
% Title block hooks (\pretitle, \posttitle, \predate, ...); load before authblk
\usepackage{titling}
\setlength{\droptitle}{-5em}

% COLORS & TEXT LAYOUT
% Named SVG colors and table row coloring
\usepackage[table, svgnames]{xcolor}
% Line spacing: \onehalfspacing, \doublespacing
\usepackage{setspace}
% Paragraphs: vertical space instead of indentation
\usepackage{parskip}
% Comment blocks: \begin{comment} ... \end{comment}
\usepackage{comment}

% LISTS
\usepackage{enumitem}
% Tighter itemize spacing
\setlist[itemize]{topsep=0pt, partopsep=0pt, parsep=0pt}

% HEADERS & FOOTERS
\usepackage{fancyhdr}
\setlength{\headheight}{14.49998pt} % Fixes fancyhdr warning with 12pt text
\fancyhf{}
% Footer: "X of Y", centered. The label mylastpage is set in config/07_macros.tex
\fancyfoot[C]{\thepage\ of \pageref{mylastpage}}
\pagestyle{fancy}
% Plain pages (title page, chapter starts): no header, no rules
\fancypagestyle{plain}{
  \fancyhf{}
  \fancyfoot[C]{\thepage\ of \pageref{mylastpage}}
  \renewcommand{\headrulewidth}{0pt}
  \renewcommand{\footrulewidth}{0pt}
}

% TITLE BLOCK
% Multiple authors via \author{...} and \affil{...}
\usepackage{authblk}
\renewcommand\Authfont{\small}
% Separate authors with commas instead of "and"
\renewcommand\Authsep{, }
\renewcommand\Authand{, }
\renewcommand\Authands{, }
            % colors, spacing, headers/footers, lists, title block
% Math environments (align, gather, ...) and symbols
\usepackage{amsmath}
% Bold math: \bm{...}
\usepackage{bm}
              % math packages
% GRAPHICS & CAPTIONS
% [H] placement specifier for exact float position
\usepackage{float}
\usepackage{graphicx}
% Keeps figures inside the margins (adjustbox)
\usepackage{adjustbox}
\graphicspath{{img/}} % Folder with image assets
% TODO: remove if no section uses it
\newcounter{figcounter}
% Captions: small font
\usepackage{caption}
\captionsetup{font=small}
% Side-by-side figures with (a), (b) sub-captions
\usepackage{subcaption}

% TIKZ & PLOTS
\usepackage{tikz}
\usepackage{pgfplots}
\pgfplotsset{compat=1.18}
\usetikzlibrary{3d}

% TABLES & CODE
% Code listings
\usepackage{listings}
% tabularray tables
\usepackage{tabularray}
% Tabularray templates: no repeated head/foot across page breaks,
% caption placed in the foot
\DefTblrTemplate{firsthead}{default}{}
\DefTblrTemplate{middlehead}{default}{}
\DefTblrTemplate{lasthead}{default}{}
\DefTblrTemplate{firstfoot}{default}{\UseTblrTemplate{caption}{default}}
\DefTblrTemplate{foot}{default}{\UseTblrTemplate{caption}{default}}
    % graphics, captions, TikZ, tables, listings
% Requires babel and csquotes to be loaded first (config/01_language.tex).
\usepackage[backend=biber, style=alphabetic, sorting=ydnt]{biblatex}
\addbibresource{references.bib} % Bibliography database
      % biblatex/biber
% Keep this file late in setup.tex. Packages that need hyperref (e.g. cleveref) must be loaded after it.

% URL line breaking (xurl also loads url)
\usepackage{xurl}
% Clickable references, citations and URLs
\usepackage[colorlinks=true, linkcolor=blue, citecolor=blue, urlcolor=blue, pdfborder={0 0 0}, hyperfootnotes=true]{hyperref}
% Links to figures/tables jump to the top of the float (needs hyperref)
\usepackage[table, figure]{hypcap}
        % url + hyperref (keep late)
% CUSTOM COMMANDS

% Clickable monospaced email: \emailaddr{name@example.com}
\newcommand{\emailaddr}[1]{\href{mailto:#1}{\texttt{#1}}}

% End of document with bibliography.
% The last-page label comes after the bibliography, so "X of Y" counts it.
\newcommand{\conclusion}{
  \clearpage
  \printbibliography
  \label{mylastpage}
}

% End of document with bibliography kept outside the page count.
% The label comes before the bibliography, and the bibliography pages have no footer.
\newcommand{\conclusionseparate}{
  \label{mylastpage}
  \clearpage
  \pagestyle{empty}
  \printbibliography
}

% End of document without bibliography (documents with no citations).
% Keeps the "X of Y" footer working.
\newcommand{\conclusionnobib}{
  \label{mylastpage}
}

% Compact title block for documents with no authors and no date:
% removes the empty author/date lines and pulls the text up under the title.
% Use in the preamble right after \title{...}
\newcommand{\compacttitle}{
  \author{}\date{}
  \predate{}\postdate{}
  \posttitle{\par\end{center}\vspace{-2em}}
}
            % custom commands


\title{\textbf{Project Assignment Case: MMORPG Company}}
\compacttitle

\begin{document}

\maketitle

\section*{Background}

This project assignment is part of the EMA and IS1 courses. You will work in a group of four or five students, and you will need
approximately ten days to complete the assignment.

The assignment concerns a start-up company that develops, operates, and sells a massively multiplayer online role-playing game, or
MMORPG. Your task is to design both the company and its product offering. This includes defining the game, its intended players, its
distinctive value, and the principal activities, IT systems, and services needed to develop and operate it.

You should approach the assignment as if you were establishing your own company. The case description below provides a common
starting point, but it is intentionally incomplete. You are expected to extend it with your own ideas and assumptions. You may also
modify some of the suggested arrangements if you clearly explain and justify your decisions.

The purpose is not to design every aspect of a real game or game company in detail. Instead, you should create a limited but
coherent company design that can be represented through the enterprise models and other deliverables required in the assignment.

\section*{The game and its distinctive offering}

The company's principal product is its MMORPG. An important part of your assignment is therefore to define what kind of game the
company will offer and why people would choose to play it.

You should develop an overall concept for the game. For example, you may consider its setting, theme, visual identity, player
activities, social features, progression system, and general playing experience. The game could focus on exploration, combat,
cooperation, competition, construction, trade, storytelling, social interaction, or some combination of these elements.

You are not expected to produce a detailed game-design document, write a complete story, or specify every game mechanic. Focus on
the small number of characteristics that define the game and distinguish it from competing offerings.

Your group should be able to explain:

\begin{itemize}
  \item what players primarily do in the game;
  \item what kind of experience the game is intended to provide;
  \item what is distinctive or innovative about the game;
  \item why players would choose it rather than another MMORPG; and
  \item what constitutes the game's unique selling proposition, or USP.
\end{itemize}

The USP should be sufficiently specific to guide the design of the company. A statement such as \enquote{the game is enjoyable and
has good graphics} is too general because the same claim could be made about almost any game. In contrast, a game might offer an
unusual form of cooperation, a distinctive virtual world, player-created societies, accessible gameplay, an innovative use of AI,
or another clearly identifiable advantage.

\section*{Players and target market}

You should define the game's principal target market. \enquote{People who play MMORPGs} is too broad to constitute a useful target
market. Instead, identify one or a few player groups whose interests, needs, and playing behaviour influence the design of the game.

For example, you may consider players' preferred game experiences, level of experience, age group, geographical market, available
playing time, social motivations, accessibility needs, and willingness to pay. You do not need to conduct extensive market research,
but your description should be plausible and sufficiently precise to support your company design.

Explain what the selected players value and how the game responds to their needs. Your target market, game concept, USP, marketing
approach, distribution channels, and revenue model should be mutually consistent.

\section*{Developing the company}

The company must possess or obtain the capabilities required to create and deliver the proposed game experience. For example, a
company offering a continuously changing virtual world may require strong capabilities in content development and technical
operations. A company built around player-created content may instead require capabilities for creator support, content review,
intellectual-property management, and community governance.

You should therefore consider not only what is attractive about the game but also how the company can reliably produce and maintain
that offering.

Brainstorming is a significant part of the project. Begin by developing ideas within your group. After generating some initial ideas
yourselves, you may use an LLM tool, such as ChatGPT or Gemini, to explore alternatives, identify weaknesses, or develop your ideas
further. The final company design should nevertheless be your group's own coherent and justified proposal.

You can obtain general background information about MMORPGs from
\href{https://en.wikipedia.org/wiki/Massively_multiplayer_online_role-playing_game}{Wikipedia} and examples of existing games from
\href{https://www.mmorpg.com/}{MMORPG.com}.

\section*{Distribution and revenues}

The company primarily provides its game to individual consumers and requires revenue to both develop the game and operate the online
service.

The initial idea is that customers pay an upfront purchase price and then a monthly subscription fee. However, your group may retain,
modify, or replace this model. For example, you may propose subscription-only access, free-to-play access with optional purchases,
different membership levels, paid expansions, or another model. Your choice should be consistent with the target market and the
experience the company intends to offer.

The game can be distributed through digital platforms such as Steam and through the company's own website. Your group may also
introduce other appropriate distribution channels, such as cloud-gaming services or physical collector's editions.

The company has agreements with several Internet service providers. An ISP may include access to the game as part of one of its
customer packages. Your group should decide what the agreement includes and how the game company and the ISP benefit from it.

\section*{Virtual products and third-party suppliers}

The company sells virtual products directly to players. These products may include armour, weapons, clothing, buildings, land areas,
decorative objects, or other items appropriate for your game.

Most products are pre-designed or configurable. However, the company will also offer a limited customisation service, allowing a
player to participate in designing a special item before deciding whether to purchase it.

The company also operates an authorised marketplace in which selected third-party suppliers can sell virtual products. Suppliers may
create the products themselves or acquire products from the company or players. Only suppliers authorised by the company may sell
through the marketplace, and the company receives a fee from each transaction.

\section*{The Game Forum}

A small group of well-known gamers intends to establish an organisation called the Game Forum. It gathers ideas from players
worldwide and helps the company identify promising improvements to the game. Particularly valuable ideas may receive financial
rewards, while selected contributors and players may receive publicity through articles, competitions, or ranking lists.

The company sponsors the Game Forum because it provides design ideas, strengthens the player community, and functions as a marketing
channel. The Game Forum may also receive revenues from advertisements, events, memberships, or other sources.

\section*{Development and operation}

An MMORPG is both a product and a continuously operating digital service. The company must therefore perform activities such as
developing and testing software, releasing content, operating technical infrastructure, supporting customers, managing the player
community, and protecting the game against cybersecurity threats and cheating.

You do not need to model all these activities. Select the processes and capabilities that are most important for the game concept,
target market, and USP proposed by your group.

\section*{Artificial intelligence}

The company is considering using artificial intelligence to improve the game experience, its internal operations, or both. Possible
applications include adaptive non-player characters, personalised quests, AI-supported creation of environments or virtual products,
customer support, content moderation, cheating detection, analysis of player behaviour, and support for software development and
testing.

Your group should select one or two meaningful AI applications. For each application, explain what need or opportunity it addresses,
who receives value from it, and what organisational and technical capabilities it requires. AI should support the company's game
concept and goals rather than be introduced merely because it is fashionable.

You should also identify one or two important risks associated with the selected applications, such as privacy problems,
inappropriate generated content, unfair automated decisions, intellectual-property concerns, cybersecurity risks, or insufficient
human oversight. Propose reasonable measures for managing these risks. A detailed legal or technical analysis is not required.

The company may develop its AI capabilities internally, obtain them from external providers, cooperate with specialist partners, or
combine these alternatives. Your enterprise architecture should reflect the approach chosen by your group.

\section*{Creating a coherent company design}

Your group has considerable freedom, but the different parts of the proposed company should be mutually consistent. The target market
should influence the game concept; the game concept and USP should influence the company's goals and processes; and the IT systems
and enterprise architecture should support the company's way of working. The selected KPIs should indicate whether the company is
succeeding in delivering its intended player experience and achieving its principal business goals.

Focus on a manageable number of well-developed ideas. A smaller and coherent company design is preferable to an extensive proposal
containing many unrelated features. Clearly state important assumptions and briefly justify your most significant design decisions.

\conclusionnobib

\end{document}
